\section{Критерий обратимости оператора}

\begin{theorem}[Критерий обратимости]
Пусть $V$ --- конечномерное линейное пространство ($\dim V = n$), и $\varphi \in L(V, V)$. Следующие условия эквивалентны:
\begin{enumerate}
    \item Оператор $\varphi$ обратим (существует $\varphi^{-1}$).
    \item Оператор $\varphi$ является биекцией (взаимно однозначным отображением) пространства $V$ на себя.
    \item Матрица $A_\varphi$ оператора $\varphi$ в любом (и, следовательно, в любом) базисе невырождена ($\det A_\varphi \neq 0$).
    \item Ядро оператора тривиально: $\ker \varphi = \{0\}$.
    \item Образ оператора совпадает со всем пространством: $\operatorname{Im} \varphi = V$.
\end{enumerate}
\end{theorem}

\begin{proof}
Докажем цепочку импликаций $(1) \Rightarrow (2) \Rightarrow (4) \Rightarrow (5) \Rightarrow (3) \Rightarrow (1)$.

\begin{enumerate}
    \item \textbf{$(1) \Rightarrow (2)$:} Если существует $\varphi^{-1}$, то по определению $\varphi \circ \varphi^{-1} = \varphi^{-1} \circ \varphi = I$. Это в точности означает, что $\varphi$ сюръективно и инъективно, то есть биективно.
    
    \item \textbf{$(2) \Rightarrow (4)$:} Пусть $\varphi$ биективно. Инъективность означает, что из $\varphi(x_1) = \varphi(x_2)$ следует $x_1 = x_2$. В частности, $\varphi(x) = 0 = \varphi(0) \Rightarrow x = 0$. Следовательно, $\ker \varphi = \{0\}$.
    
    \item \textbf{$(4) \Rightarrow (5)$:} По теореме о размерности ядра и образа: $\dim \ker \varphi + \dim \operatorname{Im} \varphi = \dim V = n$. Если $\ker \varphi = \{0\}$, то $\dim \ker \varphi = 0$, следовательно, $\dim \operatorname{Im} \varphi = n$. Поскольку $\operatorname{Im} \varphi$ --- подпространство в $V$ той же размерности, то $\operatorname{Im} \varphi = V$.
    
    \item \textbf{$(5) \Rightarrow (3)$:} Пусть $\operatorname{Im} \varphi = V$. Возьмем произвольный базис $E$. Матрица $A_\varphi$ имеет столбцы, являющиеся координатами векторов $\varphi(e_1), \dots, \varphi(e_n)$. Так как эти векторы порождают $V$, они линейно независимы (их $n$ штук в $n$-мерном пространстве). Следовательно, матрица $A_\varphi$ имеет полный ранг $n$, а значит, $\det A_\varphi \neq 0$.
    
    \item \textbf{$(3) \Rightarrow (1)$:} Пусть $\det A_\varphi \neq 0$. Тогда существует обратная матрица $A_\varphi^{-1}$. По теореме о задании оператора на базисе, существует оператор $\psi$, матрицей которого в базисе $E$ является $A_\varphi^{-1}$. По теореме о матрице произведения, матрица композиции $\varphi \circ \psi$ равна $A_\varphi A_\varphi^{-1} = E_n$, то есть совпадает с матрицей тождественного оператора $I$. В силу единственности матрицы оператора, $\varphi \circ \psi = I$. Аналогично $\psi \circ \varphi = I$. Значит, $\psi = \varphi^{-1}$, и оператор обратим.
\end{enumerate}
Все условия эквивалентны.
\hfill$\blacksquare$
\end{proof}

\begin{remark}
В конечномерном пространстве инъективность линейного оператора автоматически влечет его сюръективность, и наоборот. Это специфическое свойство, не выполняющееся в бесконечномерных пространствах.
\end{remark}
